% Sections 8.4 to 8.6, from lectures/L08/appendix/a_state_machines.md, parts A.4 to A.6: the
% machine in VHDL, Mealy machines, and the board. The two worked examples, fsm_led and
% seq_detect_mealy, are printed in full at the end of the subsections that introduce them.

\section{The Same Machine in VHDL}
\label{c8:a4}
That was the last time in this book you will design a machine by hand. Everything from here writes
it directly in VHDL, which is worth having earned rather than started with: the synthesis tool does
exactly what you just did on paper, and when a state machine misbehaves you will be debugging the
diagram, not the VHDL.

The translation replaces the explicit state table and Karnaugh equations with an \textbf{enumerated
type} naming each state, a \textbf{case statement} describing each state's transitions, and the
synthesis tool left to work out the flip-flop encoding and gate-level logic.

\subsection{The enumerated state type}
\label{c8:enum}
Suppose a machine with three states controlling an LED: \code{STATE_OFF}, \code{STATE_BLINK}
(blinking at a fixed rate) and \code{STATE_ON}. VHDL declares exactly those three values as a new
type:

\begin{vhdlcode}
type state_t is (STATE_OFF, STATE_BLINK, STATE_ON);
signal state: state_t;
...
state <= STATE_ON;
\end{vhdlcode}

\subsection{The case-statement pattern}
\label{c8:case}
The machine is almost always a synchronous process updating \code{state} on the rising edge via a
\code{case}, one branch per state, with an asynchronous reset returning it to a known starting
state. Assuming \code{to_next_state} and \code{to_prev_state} are synchronized, edge-detected pulses
produced exactly like \code{X} in \secref{c8:a2} and \secref{c8:a3}:

\begin{vhdlcode}
process(clock, reset_s2_n) is
begin
    if (reset_s2_n = '0') then
        -- On reset, go to STATE_OFF.
        state <= STATE_OFF;
    elsif (rising_edge(clock)) then
        case (state) is
            when STATE_OFF =>
                if (to_next_state = '1') then
                    state <= STATE_BLINK;
                elsif (to_prev_state = '1') then
                    state <= STATE_ON;
                end if;
            -- STATE_BLINK and STATE_ON follow the same shape.
            when others =>
                state <= STATE_OFF;
        end case;
    end if;
end process;
\end{vhdlcode}

\noindent\code{case (state) is ... end case;} must cover every value \code{state_t} can take, and
\code{when others} keeps it exhaustive if you add a state later and forget a branch.

What \code{when others} is \emph{not} is a safety net against a corrupted state register. Once
every state has its own branch, as in \file{fsm_led.vhd} below (the listing above abbreviates two
of them away, so \emph{its} \code{when others} is reachable and would wrongly reset the machine from
\code{STATE_BLINK} or \code{STATE_ON}), every value of \code{state_t} is already named above it, so
the branch is unreachable in the source and synthesis normally optimizes it away rather than
building recovery logic. Once Quartus re-encodes the machine, most bit patterns are illegal, and
recovering from one requires Quartus's Safe State Machine setting rather than anything you can
write in VHDL.

\subsection{Worked example: \code{fsm_led}}
\label{c8:fsmled}
\file{lectures/L08/fsm_led/fsm_led.vhd} implements exactly that three-state machine, extended to two
buttons: \code{button_n(0)} advances, \code{button_n(1)} goes back, the states are closed in a loop,
and \code{STATE_BLINK} toggles the LED every 100\,ms.

Three subcomponents are reused, all of them modules you wrote yourself:
\begin{itemize}
  \item \code{reset_sync}: Chapter~\ref{c4:ch}'s reset synchronizer, unchanged
    (Exercise~\ref{c4:ex:8}).
  \item \code{button_sync}: Chapter~\ref{c4:ch}'s button synchronizer, also unchanged, but
    instantiated with \code{generic map(2)} so it handles both buttons at once, producing a two-bit
    pulse vector. This is the payoff of the generic: the file is identical to the one
    Chapters~\ref{c4:ch} and~\ref{c7:ch} instantiate with \code{1} (Exercise~\ref{c4:ex:9}).
  \item \code{timer}: Chapter~\ref{c7:ch}'s timer, unchanged (Exercise~\ref{c7:ex:4}).
    \code{fsm_led} declares its own \code{TIMER_TICK_COUNT} generic and passes it straight through,
    defaulting to \code{5_000_000} for 100\,ms at \code{50 MHz}, and only enables the timer while in
    \code{STATE_BLINK}.
\end{itemize}
None of the three is shipped in \file{fsm_led/}: copy your own in before building the example. By
this point in the book that is four chapters' worth of your own modules composed into one design,
which is the thing worth noticing about it.

Two synchronous processes implement the machine:
\begin{itemize}
  \item \code{STATE_PROCESS} is the \code{case} transition logic above, driven by
    \code{to_next_state}/\allowbreak\code{to_prev_state}, derived combinationally from
    \code{button_edge_s2}.
    Each requires the \emph{other} button's pulse to be \code{'0'} in the same cycle, so pressing
    both close enough together that their pulses land on one clock cycle moves the machine nowhere
    rather than picking a winner arbitrarily. An ambiguous input deserves a defined answer, and
    ``do nothing'' is the one that cannot surprise you.
  \item \code{LED_PROCESS} drives \code{led_s} purely from the current state: off in
    \code{STATE_OFF}, on in \code{STATE_ON}, toggling on every timer timeout in \code{STATE_BLINK}.
    Because it never reads \code{button_n} or \code{button_edge_s2} directly, only \code{state},
    this is a genuine Moore machine: the LED reacts to a press only indirectly, through the state it
    caused.
\end{itemize}
This is the whole of \file{lectures/L08/fsm_led/fsm_led.vhd}:

\vhdlfile{lectures/L08/fsm_led/fsm_led.vhd}

\section{Mealy Machines: The One-Cycle Difference}
\label{c8:a5}
Everything so far has been a \textbf{Moore} machine, deliberately: Moore is the safer default and
the form you design by hand. A \textbf{Mealy} machine relaxes the rule that defines a Moore machine:
the output may depend on the current state \emph{and} the current input, computed combinationally.
It is not stored in a register, and it changes within the same clock cycle the input does.

\subsection{The behaviour: detect two consecutive \code{1} bits}
\label{c8:mealydetect}
A serial input \code{din} carries one bit per clock cycle, and \code{y} should indicate that the
last two bits were both \code{1}, with overlapping matches counting, so \code{111} reports a match
twice.

As a Mealy machine this needs only two states, because the output can react to the input
\emph{before} it is reflected in the next state: \code{STATE_IDLE} (the previous bit was \code{0},
or this is the first) and \code{STATE_ONE} (the previous bit was \code{1}).

\begin{narrowtable}{@{}lcll@{}}
  \thead{State} & \thead{\code{din}} & \thead{Next state} & \thead{\code{y}} \\
  \midrule
  \code{STATE_IDLE} & 0 & \code{STATE_IDLE} & 0 \\
  \code{STATE_IDLE} & 1 & \code{STATE_ONE} & 0 \\
  \code{STATE_ONE} & 0 & \code{STATE_IDLE} & 0 \\
  \code{STATE_ONE} & 1 & \code{STATE_ONE} & \textbf{1} \\
\end{narrowtable}

\noindent The state process is \secref{c8:a4}'s \code{case} pattern; only the output differs. This
is \file{lectures/L08/seq_detect_mealy/seq_detect_mealy.vhd}, and its output assignment is a single
concurrent statement, entirely outside any clocked process:

\begin{vhdlcode}
y <= '1' when (state = STATE_ONE and din = '1') else '0';
\end{vhdlcode}

\noindent That is a \textbf{conditional signal assignment}, and it is the last piece of VHDL syntax
this book introduces. It is a concurrent statement, so it lives in the architecture body rather than
inside a process, and it describes a multiplexer: pick the first value whose condition holds,
otherwise the value after \code{else}. It is the same hardware you would get from an
\code{if}/\code{else} inside a combinational process, written in one line instead of five, and it is
worth having because a one-line output assignment makes it obvious at a glance what the output
depends on.

\begin{vhdlcode}
-- These two describe exactly the same multiplexer.
x <= a when sel = '1' else b;

process(sel, a, b) is
begin
    if (sel = '1') then
        x <= a;
    else
        x <= b;
    end if;
end process;
\end{vhdlcode}

\noindent Always give it an \code{else}. Leave one out and you have described a signal that keeps
its old value when no condition holds, which is a latch, exactly as in \secref{c5:a3}.
\code{fsm_led} uses this form three times, and Exercise~\ref{c8:ex:4} asks you to use it.

\textbf{That one line is the whole distinction.} It reads \code{din}, so it is Mealy. Move the same
condition inside the clocked process, or drop \code{din} from it, and you have a Moore machine
instead.

This is the whole of \file{lectures/L08/seq_detect_mealy/seq_detect_mealy.vhd}:

\vhdlfile{lectures/L08/seq_detect_mealy/seq_detect_mealy.vhd}

\subsection{What it costs, and what it buys}
\label{c8:mealycost}
The Moore version of this detector needs a \textbf{third} state, \code{M_TWO}, existing purely to
remember for one cycle that the match happened, so the output has something to read that does not
depend on the input. Feed both machines \code{1, 1, 0, 1, 1, 1}:

\begin{narrowtable}{@{}cclclc@{}}
  \thead{Cycle} & \thead{\code{din}} & \thead{Mealy state (before)} & \thead{Mealy \code{y}} &
    \thead{Moore state (before)} & \thead{Moore \code{y}} \\
  \midrule
  1 & 1 & \code{STATE_IDLE} & 0 & \code{M_IDLE} & 0 \\
  2 & 1 & \code{STATE_ONE} & \textbf{1} & \code{M_ONE} & 0 \\
  3 & 0 & \code{STATE_ONE} & 0 & \code{M_TWO} & \textbf{1} \\
  4 & 1 & \code{STATE_IDLE} & 0 & \code{M_IDLE} & 0 \\
  5 & 1 & \code{STATE_ONE} & \textbf{1} & \code{M_ONE} & 0 \\
  6 & 1 & \code{STATE_ONE} & \textbf{1} & \code{M_TWO} & \textbf{1} \\
\end{narrowtable}

\noindent The Moore \code{y} column is exactly the Mealy \code{y} column shifted down one row. That
is the entire trade: Mealy reports the match the \emph{same} cycle the second \code{1} arrives where
Moore takes one more, and Mealy needs 2 states where Moore needs 3.

\subsection{Which to reach for}
\label{c8:whichone}
Moore, almost always, and it is what the rest of this book uses:
\begin{itemize}
  \item its output depends on no primary input, so it cannot follow a noisy one.
  \item a Mealy output can in principle glitch combinationally if its inputs aren't already clean.
  \item note that ``Moore'' on its own does not mean glitch-free. An output decoded combinationally
    from several state bits, like \secref{c8:a2}'s \code{Y = Q1Q2'}, can still glitch while those
    bits skew on a transition. What makes an output glitch-free is \emph{registering} it, the way
    \code{fsm_led}'s \code{LED_PROCESS} does, and that is what you want on anything driving a wire
    off the chip. A serial \code{tx} line decoded straight from the state bits will emit spurious
    edges.
  \item a machine whose output depends on fewer things is a machine with fewer ways to be wrong.
\end{itemize}
Reach for Mealy only when you specifically need the lower latency or the smaller state count, and
know that you are trading a synchronous output for a combinational one when you do.

\section{FPGA Demonstration}
\label{c8:a6}
\code{fsm_led} synthesizes and runs exactly the way every VHDL design has since
Chapter~\ref{c1:ch} (Appendix~\ref{app:quartus}): create a Quartus Prime Lite project targeting the
DE0-CV's \code{5CEBA4F23C7N}, add the \file{.vhd} files for the example and its subcomponents,
compile, assign each port to a physical pin in the Pin Planner, then recompile and program the
board.

For \code{fsm_led}, assign \code{clock} to the 50\,MHz oscillator pin, \code{reset_n} and
\code{button_n(1 downto 0)} to three push buttons, and \code{led} to an onboard LED. Then confirm
that pressing ``next'' repeatedly cycles the LED \code{off -> blinking -> on -> off}, that
``previous'' runs the other way, and that reset returns it to off.

\code{seq_detect_mealy} is deliberately \textbf{not} demonstrated on the board. It consumes one bit
per clock cycle, and a switch flipped by hand holds its value for millions of cycles, so the only
thing an LED can show is the degenerate case: \code{din} held high, \code{y} high from the second
cycle onward. The one-cycle distinction the whole of \secref{c8:a5} is about happens far too fast
to see. Run its testbench instead, which is the general lesson rather than an exception:

\begin{shell}
cd lectures/L08/seq_detect_mealy
cp ../../L04/exercises/reset_sync/reset_sync.vhd .     # the one you wrote in L04
ghdl -a --std=93 reset_sync.vhd seq_detect_mealy.vhd seq_detect_mealy_tb.vhd
ghdl -e --std=93 seq_detect_mealy_tb
ghdl -r --std=93 seq_detect_mealy_tb --assert-level=error --stop-time=10ms
\end{shell}

\noindent Exact pin numbers depend on which DE0-CV buttons and LEDs you choose, assigned the same way
as for every design since Chapter~\ref{c1:ch}. Refer to the DE0-CV's own documentation for its
physical pinout if you need to confirm which header pin corresponds to which labelled button, switch
or LED.
